\documentclass[11pt,a4paper]{article}

\usepackage[T2A]{fontenc}
\usepackage[utf8]{inputenc}
\usepackage[english,russian]{babel}
\usepackage{amsmath,amssymb,amsthm,mathtools}
\usepackage{geometry}
\geometry{left=2.2cm,right=2.0cm,top=2.0cm,bottom=2.2cm}
\usepackage{enumitem}
\usepackage{booktabs}
\usepackage{xcolor}
\usepackage[unicode,hidelinks]{hyperref}

\renewcommand{\proofname}{Доказательство}
\newtheorem{thm}{Теорема}[section]
\newtheorem{lem}[thm]{Лемма}
\newtheorem{utv}[thm]{Утверждение}
\theoremstyle{definition}
\newtheorem{defn}[thm]{Определение}
\newtheorem{zam}[thm]{Замечание}
\newtheorem{prim}[thm]{Пример}

\newcommand{\PP}{\mathbf P}
\newcommand{\EE}{\mathbf E}
\newcommand{\R}{\mathbb R}
\newcommand{\N}{\mathbb N}
\newcommand{\Q}{\mathbb Q}
\newcommand{\B}{\mathcal B}
\newcommand{\F}{\mathcal F}
\newcommand{\1}{\mathbf 1}
\newcommand{\eqd}{\overset{d}{=}}

\newenvironment{task}[1]{\par\medskip\noindent\colorbox{black!7}{\parbox{\dimexpr\linewidth-2\fboxsep}{\textbf{#1}}}\par\smallskip}{\par}
\newenvironment{answer}{\par\noindent\textbf{Ответ.}~}{\par}

\title{\textbf{Случайные процессы. Листок SP-1}\\[4pt]
\large Теорминимум и решения задач\\[2pt]
\normalsize Борелевская $\sigma$-алгебра, случайные последовательности, цепи Маркова}
\date{}

\begin{document}
\maketitle
\tableofcontents
\newpage

\part*{Часть I. Теорминимум}
\addcontentsline{toc}{part}{Часть I. Теорминимум}

Этот раздел написан в расчёте на то, что теория вероятностей подзабыта. Сначала
коротко повторяется всё нужное из базового курса, затем по порядку идут темы
листка Б-0 -- Б-3. Всё, что используется в решениях, здесь есть.

\section{Что нужно вспомнить из теории вероятностей}

\subsection{Вероятностное пространство}

Вероятностное пространство~--- тройка $(\Omega,\F,\PP)$:
\begin{itemize}[nosep]
\item $\Omega$~--- множество элементарных исходов;
\item $\F$~--- \emph{$\sigma$-алгебра} событий: семейство подмножеств $\Omega$,
содержащее $\Omega$, замкнутое относительно дополнения и \emph{счётных}
объединений (а значит, и счётных пересечений);
\item $\PP$~--- вероятностная мера: $\PP(\Omega)=1$ и
$\PP\bigl(\bigsqcup_n A_n\bigr)=\sum_n\PP(A_n)$ для попарно непересекающихся $A_n$.
\end{itemize}
Главное слово здесь~--- \textbf{счётных}. Несчётные объединения событий событиями
быть не обязаны. Поэтому во всех задачах на измеримость приходится сводить
«для всех $\varepsilon>0$» к «для всех $\varepsilon=1/m$», а «существует
действительное $c$»~--- к «существует рациональное $q$».

Полезные следствия аксиом: непрерывность меры~--- если $A_n\uparrow A$, то
$\PP(A_n)\to\PP(A)$; если $A_n\downarrow A$, то тоже $\PP(A_n)\to\PP(A)$.

\subsection{Случайная величина и её распределение}

\emph{Случайная величина}~--- функция $X:\Omega\to\R$, для которой
$\{X\le x\}=\{\omega: X(\omega)\le x\}\in\F$ при всех $x\in\R$ (эквивалентно:
$\{X\in B\}\in\F$ для всех борелевских $B$). Это и есть «измеримость».
\emph{Распределение} $X$~--- мера $\PP_X(B)=\PP(X\in B)$ на $\R$. Функция
распределения $F_X(x)=\PP(X\le x)$ однозначно определяет распределение.

Если $X$ дискретна, распределение задаётся числами $\PP(X=x_k)$; если есть
плотность $p_X$, то $\PP(X\in B)=\int_Bp_X(x)\,dx$.

Суммы, произведения, $\max$, $\min$, $\sup_n$, $\inf_n$, $\limsup_n$, $\liminf_n$
счётного числа случайных величин~--- снова случайные величины (возможно, со
значениями $\pm\infty$). Непрерывная (и вообще борелевская) функция от
случайной величины~--- случайная величина.

\subsection{Независимость, условная вероятность}

События $A,B$ независимы, если $\PP(AB)=\PP(A)\PP(B)$. Величины
$X_1,\dots,X_n$ независимы, если
$\PP(X_1\in B_1,\dots,X_n\in B_n)=\prod\PP(X_k\in B_k)$ для всех борелевских
$B_k$. Функции от непересекающихся групп независимых величин независимы.

Условная вероятность: $\PP(A\mid B)=\PP(AB)/\PP(B)$ при $\PP(B)>0$.
Формула умножения:
\[
\PP(A_1A_2\dots A_n)=\PP(A_1)\PP(A_2\mid A_1)\PP(A_3\mid A_1A_2)\cdots
\PP(A_n\mid A_1\dots A_{n-1}).
\]
Формула полной вероятности: если $\Omega=\bigsqcup_kH_k$, то
$\PP(A)=\sum_k\PP(A\mid H_k)\PP(H_k)$.

Для величин с плотностями условная плотность
$p_{Y\mid X}(y\mid x)=p_{X,Y}(x,y)/p_X(x)$.

\subsection{Характеристическая функция}

\[
\varphi_X(t)=\EE e^{itX},\qquad
\varphi_{(X_1,\dots,X_k)}(t_1,\dots,t_k)=\EE e^{i(t_1X_1+\dots+t_kX_k)} .
\]
Свойства, которые нужны:
\begin{enumerate}[nosep]
\item \textbf{теорема единственности}: х.ф. однозначно определяет распределение
(и для величин, и для векторов);
\item $\varphi_{aX+b}(t)=e^{itb}\varphi_X(at)$;
\item если $X,Y$ независимы, то $\varphi_{X+Y}=\varphi_X\varphi_Y$;
\item х.ф. проекции: $\varphi_{(X_1,\dots,X_{k})}(t_1,\dots,t_k)=
\varphi_{(X_1,\dots,X_{k+1})}(t_1,\dots,t_k,0)$ (подставили ноль~--- «забыли»
координату).
\end{enumerate}

\subsection{Гауссовские величины и векторы}

$X\sim\mathcal N(a,\sigma^2)$: плотность
$\frac{1}{\sqrt{2\pi}\sigma}e^{-(x-a)^2/(2\sigma^2)}$, х.ф.
$e^{ita-\sigma^2t^2/2}$. Вектор называется гауссовским, если любая линейная
комбинация его координат~--- гауссовская величина. Распределение гауссовского
вектора полностью задаётся средними и ковариационной матрицей. Линейное
преобразование гауссовского вектора~--- снова гауссовский вектор. Независимые
гауссовские величины образуют гауссовский вектор.

\subsection{Теорема о единственности меры}

\begin{thm}[о $\pi$-системах]
Если две вероятностные меры совпадают на семействе множеств $\mathcal C$,
замкнутом относительно конечных пересечений ($\pi$-системе), то они совпадают и
на $\sigma(\mathcal C)$~--- наименьшей $\sigma$-алгебре, содержащей $\mathcal C$.
\end{thm}
Именно поэтому функция распределения определяет распределение: множества
$(-\infty,x]$ образуют $\pi$-систему, порождающую $\B(\R)$.

\section{Б-0. Борелевская $\sigma$-алгебра}

\subsection{Определения}

\begin{defn}
$\B(\R)$~--- наименьшая $\sigma$-алгебра, содержащая все интервалы (или все лучи
$(-\infty,x]$~--- получится то же самое). $\B(\R^n)$~--- наименьшая
$\sigma$-алгебра, содержащая все «параллелепипеды» $B_1\times\dots\times B_n$.
\end{defn}

Пространство последовательностей $\R^\infty=\{x=(x_1,x_2,\dots)\}$.
\emph{Цилиндром} называется множество вида
\[
C=\{x\in\R^\infty:\ (x_1,\dots,x_n)\in B\},\qquad B\in\B(\R^n),
\]
то есть условие наложено лишь на \emph{конечное} число координат.

\begin{defn}
$\B(\R^\infty)$~--- наименьшая $\sigma$-алгебра, содержащая все цилиндры.
Равносильно: наименьшая $\sigma$-алгебра, относительно которой измеримы все
координатные отображения $x\mapsto x_n$.
\end{defn}

\subsection{Как доказывать, что множество борелевское}

Основной приём~--- переписать условие через кванторы и заменить кванторы на
операции над множествами. Словарь:
\begin{center}
\begin{tabular}{ll}
\toprule
Логика & Множества \\
\midrule
«и» & $\cap$ \\
«или» & $\cup$ \\
«не» & дополнение \\
$\forall n\in\N$ & $\bigcap_n$ \\
$\exists n\in\N$ & $\bigcup_n$ \\
$\forall\varepsilon>0$ & $\bigcap_{m\ge1}$ с $\varepsilon=1/m$ \\
$\exists c\in\R$ & $\bigcup_{q\in\Q}$ (если условие позволяет заменить $c$ на рациональное) \\
\bottomrule
\end{tabular}
\end{center}
«Кирпичики»~--- множества $\{x_n<c\}$, $\{x_n\le c\}$, $\{x_n>c\}$,
$\{x_n\ge c\}$, $\{x_n=c\}$: все они цилиндры. Множества вида
$\{x_n\le x_k\}$, $\{x_n=x_k\}$ тоже борелевские, потому что $x_n-x_k$~---
измеримая функция (или явно:
$\{x_n>x_k\}=\bigcup_{q\in\Q}\{x_k<q\}\cap\{x_n>q\}$).

Второй приём~--- сослаться на измеримость функций: $\sup_nx_n$, $\inf_nx_n$,
$\limsup x_n=\inf_n\sup_{k\ge n}x_k$, $\liminf x_n$ измеримы (как функции со
значениями в $[-\infty,+\infty]$), поэтому, например,
$\{\limsup x_n\le3\}$~--- прообраз борелевского множества, то есть борелевское.
На зачёте лучше уметь оба способа.

\subsection{Случайные моменты}

\begin{defn}
Пусть $X_1,X_2,\dots$~--- случайные величины. Величина $\tau$ со значениями в
$\{1,2,\dots\}\cup\{\infty\}$ называется \emph{случайным моментом}, если
$\{\tau=n\}\in\F$ при всех $n$. Если при этом $\{\tau=n\}\in\sigma(X_1,\dots,X_n)$
(«решение, остановиться ли в момент $n$, принимается по $X_1,\dots,X_n$»), то
$\tau$ называется \emph{моментом остановки} (марковским моментом).
\end{defn}

Здесь $\sigma(X_1,\dots,X_n)$~--- $\sigma$-алгебра, порождённая величинами,
то есть все события вида $\{(X_1,\dots,X_n)\in B\}$, $B\in\B(\R^n)$.
Так как $\{\tau=\infty\}=\Omega\setminus\bigcup_n\{\tau=n\}$, проверять надо
только конечные $n$.

Удобный критерий: $\tau$~--- случайный момент, если $\{\tau\le n\}\in\F$ для
всех $n$ (тогда $\{\tau=n\}=\{\tau\le n\}\setminus\{\tau\le n-1\}$).

\section{Б-1. Случайные последовательности}

\subsection{Определение и распределение}

\begin{defn}
Случайная последовательность~--- это набор случайных величин
$X=(X_1,X_2,\dots)$ на одном вероятностном пространстве. Равносильно: измеримое
отображение $X:(\Omega,\F)\to(\R^\infty,\B(\R^\infty))$.
\end{defn}
(Равносильность: прообраз цилиндра~--- событие
$\{(X_1,\dots,X_n)\in B\}\in\F$, а прообразы множеств, для которых это верно,
образуют $\sigma$-алгебру, содержащую цилиндры.)

\emph{Распределение} последовательности~--- мера
$\PP_X(A)=\PP(X\in A)$ на $\B(\R^\infty)$.
\emph{Конечномерные распределения} (к.м.р.)~--- распределения векторов
$(X_{t_1},\dots,X_{t_k})$ для всех конечных наборов индексов.

\begin{thm}
Конечномерные распределения однозначно определяют распределение
последовательности.
\end{thm}
\begin{proof}
Цилиндры образуют $\pi$-систему (пересечение двух цилиндров~--- цилиндр),
порождающую $\B(\R^\infty)$. Меры двух последовательностей с одинаковыми
к.м.р. совпадают на цилиндрах, значит, по теореме о $\pi$-системах совпадают
всюду.
\end{proof}
Для последовательности достаточно знать распределения начальных отрезков
$(X_1,\dots,X_k)$, $k=1,2,\dots$: распределение любого набора
$(X_{t_1},\dots,X_{t_k})$ получается из распределения $(X_1,\dots,X_{\max t_j})$
проекцией.

\subsection{Согласованность и теорема Колмогорова}

Пусть для каждого $k$ задана мера $P_k$ на $\R^k$ и мы хотим, чтобы это были
распределения $(X_1,\dots,X_k)$ некоторой последовательности. Очевидное
необходимое условие:
\[
\text{(С)}\qquad P_{k+1}(B\times\R)=P_k(B)\quad\text{для всех }B\in\B(\R^k),\ k\ge1
\]
(вероятность события про первые $k$ координат не зависит от того, в каком
векторе его считать). Если к.м.р. задаются для произвольных наборов индексов
$(t_1,\dots,t_k)$, добавляется ещё условие симметрии: при перестановке индексов
мера переставляется так же,
$P_{t_{\sigma(1)}\dots t_{\sigma(k)}}(B_{\sigma(1)}\times\dots\times B_{\sigma(k)})
=P_{t_1\dots t_k}(B_1\times\dots\times B_k)$, а условие (С) формулируется
как «вычёркивание» любой координаты.

\begin{thm}[Колмогорова о продолжении]
Если меры $P_k$ согласованы (условие (С)), то существует (единственная)
вероятностная мера $P$ на $(\R^\infty,\B(\R^\infty))$ с этими к.м.р., а значит,
и случайная последовательность с такими к.м.р.
\end{thm}

То же условие в других терминах (доказательства~--- в решениях базовых задач):
\begin{itemize}[nosep]
\item через функции распределения:
$\lim\limits_{x_{k+1}\to+\infty}F_{k+1}(x_1,\dots,x_k,x_{k+1})=F_k(x_1,\dots,x_k)$;
\item через характеристические функции:
$\varphi_{k+1}(\lambda_1,\dots,\lambda_k,0)=\varphi_k(\lambda_1,\dots,\lambda_k)$.
\end{itemize}

\subsection{Приём Крамера -- Уолда}

Х.ф. вектора в точке $(\lambda_1,\dots,\lambda_k)$~--- это х.ф. одной величины
$\lambda_1X_1+\dots+\lambda_kX_k$ в точке $1$:
\[
\varphi_{(X_1,\dots,X_k)}(\lambda_1,\dots,\lambda_k)=\EE e^{i(\lambda_1X_1+\dots+\lambda_kX_k)}
=\varphi_{\lambda_1X_1+\dots+\lambda_kX_k}(1).
\]
Поэтому распределения всех линейных комбинаций определяют распределение вектора.

\section{Б-2. Определение цепи Маркова}

\subsection{Дискретное пространство состояний}

Пусть $X_0,X_1,\dots$ принимают значения в конечном или счётном множестве $S$
(состояния).

\begin{defn}
$X_n$~--- \emph{цепь Маркова} (ЦМ), если для всех $n$ и всех состояний
\[
\PP(X_{n+1}=j\mid X_n=i,\,X_{n-1}=i_{n-1},\dots,X_0=i_0)=\PP(X_{n+1}=j\mid X_n=i)
\]
(всякий раз, когда условие имеет положительную вероятность). Словами: при
известном настоящем будущее не зависит от прошлого.
\end{defn}

Числа $p_{ij}(n)=\PP(X_{n+1}=j\mid X_n=i)$ образуют \emph{матрицу вероятностей
перехода} (м.в.п.) $P_n$ на шаге $n\to n+1$. Это \emph{стохастическая}
матрица: элементы неотрицательны, сумма в каждой строке равна $1$. Если $P_n=P$
не зависит от $n$, цепь \emph{однородная}. \emph{Начальное распределение}~---
вектор-строка $p^0$, $p^0_i=\PP(X_0=i)$.

\begin{thm}[совместное распределение]
\[
\PP(X_0=i_0,X_1=i_1,\dots,X_n=i_n)=p^0_{i_0}\,p_{i_0i_1}(0)\,p_{i_1i_2}(1)\cdots p_{i_{n-1}i_n}(n-1).
\]
Обратно, если совместные вероятности имеют такой вид, то $X_n$~--- ЦМ с м.в.п.
$P_n$.
\end{thm}
\begin{proof}
Первое~--- формула умножения плюс марковское свойство. Обратно: условная
вероятность
$\PP(X_{n+1}=j\mid X_n=i,\dots,X_0=i_0)$ есть отношение двух таких
произведений, и всё сокращается до $p_{ij}(n)$.
\end{proof}

Суммируя по промежуточным состояниям, получаем \textbf{уравнение
Колмогорова -- Чепмена}: вероятности перехода за несколько шагов даются
произведением матриц,
\[
\PP(X_t=j\mid X_s=i)=\bigl(P_sP_{s+1}\cdots P_{t-1}\bigr)_{ij},
\qquad\text{в однородном случае }(P^{t-s})_{ij},
\]
а распределение в момент $n$~--- это $p^n=p^0P_0P_1\cdots P_{n-1}$ (в однородном
случае $p^n=p^0P^n$).

\subsection{Главная рабочая лемма}

\begin{lem}\label{lem:rec}
Пусть $X_{n+1}=f_n(X_n,\xi_{n+1})$, где $\xi_1,\xi_2,\dots$ независимы между
собой и не зависят от $X_0$. Тогда $X_n$~--- ЦМ с
$p_{ij}(n)=\PP(f_n(i,\xi_{n+1})=j)$. Если $f_n=f$ и $\xi_n$ одинаково
распределены, цепь однородна.
\end{lem}
\begin{proof}
$(X_0,\dots,X_n)$~--- функция от $(X_0,\xi_1,\dots,\xi_n)$, а $\xi_{n+1}$ от
этого набора не зависит. Поэтому
\[
\PP(X_{n+1}=j\mid X_n=i,\dots,X_0=i_0)=
\frac{\PP(f_n(i,\xi_{n+1})=j,\ X_n=i,\dots,X_0=i_0)}{\PP(X_n=i,\dots,X_0=i_0)}
=\PP(f_n(i,\xi_{n+1})=j).
\]
\end{proof}
Почти все задачи «является ли $X_n$ цепью Маркова» с ответом «да» решаются так:
находим рекуррентное представление «новое состояние = функция от старого и
свежего независимого шума».

\subsection{Как доказывать, что процесс \emph{не} ЦМ}

Надо найти момент $n$, состояние $i$ и две предыстории
$H_1=\{X_n=i,\ \text{прошлое}_1\}$, $H_2=\{X_n=i,\ \text{прошлое}_2\}$
положительной вероятности, для которых
$\PP(X_{n+1}=j\mid H_1)\ne\PP(X_{n+1}=j\mid H_2)$. Тогда эти вероятности не
могут обе равняться $\PP(X_{n+1}=j\mid X_n=i)$. Важно: обе предыстории должны
относиться \emph{к одному и тому же моменту} $n$~--- иначе мы опровергнем
только однородность, а не марковость.

\subsection{Общее пространство состояний}

Для величин с непрерывными распределениями марковость формулируется через
условные распределения:
$\PP(X_{n+1}\in B\mid X_0,\dots,X_n)=\PP(X_{n+1}\in B\mid X_n)=P_n(X_n,B)$.
Функция $P_n(x,B)$ называется \emph{переходным ядром}; если
$P_n(x,B)=\int_Bf_n(y\mid x)\,dy$, то $f_n(y\mid x)$~--- \emph{переходная
плотность}. \emph{Переходный оператор} действует на функции:
$(P_nh)(x)=\EE[h(X_{n+1})\mid X_n=x]=\int h(y)f_n(y\mid x)\,dy$.
Для величин с плотностями марковость означает, что условная плотность
$X_{n+1}$ при условии $(X_0,\dots,X_n)$ зависит только от $X_n$; совместная
плотность тогда равна $p_0(x_0)f_0(x_1\mid x_0)\cdots f_{n-1}(x_n\mid x_{n-1})$.
Лемма~\ref{lem:rec} верна и здесь.

\section{Б-3. Свойства цепей Маркова}

\subsection{Функции от цепи}

\begin{itemize}
\item Если $h$ \emph{инъективна}, то $h(X_n)$~--- ЦМ: события
$\{h(X_n)=h(i)\}$ и $\{X_n=i\}$ совпадают, условные вероятности те же.
\item Если $h$ склеивает состояния, $h(X_n)$ может перестать быть ЦМ: склейка
теряет информацию, которая нужна для предсказания будущего (пример~--- в
решениях).
\item \textbf{Расширение состояния.} Процесс, который сам не ЦМ, часто
становится ЦМ, если добавить в состояние недостающую информацию: пара
$(X_n,X_{n+1})$, пара «позиция, направление», пара «очередь, прошедшее
время» и т.п.
\item \textbf{Прореживание.} $Y_n=X_{kn}$~--- ЦМ с м.в.п. $P^k$ (в однородном
случае).
\end{itemize}

\subsection{Строго марковское свойство}

\begin{thm}
Если $N$~--- момент остановки ($\{N=n\}\in\sigma(X_0,\dots,X_n)$) и $N<\infty$,
то при условии $X_N=i$ и всего прошлого до момента $N$ процесс
$X_{N},X_{N+1},\dots$ ведёт себя как цепь с той же м.в.п., стартующая из $i$;
в частности, $\PP(X_{N+1}=j\mid X_N=i,\ \text{прошлое})=p_{ij}$.
\end{thm}
Для случайного момента, который «заглядывает в будущее» (например, момент
последнего посещения), это неверно.

\subsection{Стационарность}

\begin{defn}
Процесс \emph{стационарен в узком смысле}, если
$(X_{n},X_{n+1},\dots,X_{n+k})\eqd(X_0,\dots,X_k)$ для всех $n,k$.
Распределение $\pi$ называется \emph{инвариантным} (стационарным) для
м.в.п. $P$, если $\pi P=\pi$ (т.е. $\pi$~--- левый собственный вектор $P$ с
собственным значением $1$, или, что то же, $P^{\mathsf T}\pi^{\mathsf T}=\pi^{\mathsf T}$).
\end{defn}

\begin{thm}
Однородная ЦМ стационарна в узком смысле $\iff$ её начальное распределение
инвариантно: $p^0P=p^0$. Для цепи с переходной плотностью условие
принимает вид $p_0(y)=\int f(y\mid x)p_0(x)\,dx$.
\end{thm}

\subsection{Обращение времени}

Если $X_n$~--- стационарная ЦМ с инвариантным распределением $\pi$, $\pi_i>0$,
то обращённый процесс $Y_n=X_{N-n}$ тоже ЦМ с м.в.п.
\[
q_{ij}=\frac{\pi_jp_{ji}}{\pi_i}.
\]
Цепь называется \emph{обратимой}, если $q_{ij}=p_{ij}$, т.е. выполнено
\emph{условие детального баланса} $\pi_ip_{ij}=\pi_jp_{ji}$ для всех $i,j$.

\newpage
\part*{Часть II. Решения}
\addcontentsline{toc}{part}{Часть II. Решения}

В группах много повторяющихся задач, поэтому каждая задача решена один раз.
Таблица соответствия:

\begin{center}
\begin{tabular}{lll}
\toprule
Группа & Номер в листке & Где решение \\
\midrule
все & нулевые 1, 2, 3 & \S\ref{sec:zero} \\
1 & 1 & задача \ref{z:cramer} \\
1 & 2, Д2 группы 3 & задача \ref{z:worker} \\
1 & 3; 3 группы 3 & задача \ref{z:inhom} \\
1 & 4; Д1 группы 3 & задача \ref{z:stop} \\
1 & Д1; 4 группы 3 & задача \ref{z:func} \\
1 & Д2; Д1 группы 2 & задача \ref{z:polygon} \\
2, 3 & 1 & задача \ref{z:consist} \\
2 & 2 & задача \ref{z:subseq} \\
2 & 3 & задача \ref{z:hom} \\
2 & 4 & задача \ref{z:station} \\
2 & Д2 & задача \ref{z:dice} \\
3 & 2 & задача \ref{z:walk} \\
--- & Ф-1.1 -- Ф-1.5 & \S\ref{sec:fac} \\
\bottomrule
\end{tabular}
\end{center}

\section{Нулевые задачи}\label{sec:zero}

Всюду $x=(x_1,x_2,\dots)\in\R^\infty$, «борелевское» означает
«принадлежит $\B(\R^\infty)$». Напомним, что $\{x:x_n<c\}$ и подобные~---
цилиндры, то есть борелевские множества.

\begin{task}{Группа 1, а) $\{x:\ \limsup x_n\le3\}$ --- борелевское.}
\end{task}
Распишем условие. $\limsup x_n\le3$ означает: для любого $\varepsilon>0$ лишь
конечное число членов больше или равно $3+\varepsilon$, т.е.
\[
\forall m\ \exists N\ \forall n\ge N:\quad x_n<3+\tfrac1m .
\]
(Проверка: если $\limsup x_n=L\le3$, то $\sup_{k\ge N}x_k\downarrow L$, поэтому
для каждого $m$ найдётся $N$ с $\sup_{k\ge N}x_k<3+\frac1m$. Обратно, если так
для всех $m$, то $\limsup x_n\le3+\frac1m$ для всех $m$, т.е. $\le3$.)
По словарю:
\[
\{\limsup x_n\le3\}=\bigcap_{m=1}^\infty\ \bigcup_{N=1}^\infty\ \bigcap_{n\ge N}\{x_n<3+\tfrac1m\}
\]
--- счётные операции над цилиндрами, значит, множество борелевское.

\begin{task}{Группа 1, б) множество монотонных последовательностей --- борелевское.}
\end{task}
Монотонная~--- это неубывающая или невозрастающая:
\[
M=\Bigl(\bigcap_{n}\{x_n\le x_{n+1}\}\Bigr)\cup\Bigl(\bigcap_{n}\{x_n\ge x_{n+1}\}\Bigr).
\]
Осталось показать, что $\{x_n\le x_{n+1}\}$ борелевское. Его дополнение
\[
\{x_n>x_{n+1}\}=\bigcup_{q\in\Q}\bigl(\{x_{n+1}<q\}\cap\{x_n>q\}\bigr)
\]
(между двумя различными числами всегда есть рациональное)~--- счётное
объединение цилиндров. Значит, $M$ борелевское.

\begin{task}{Группа 2, а) $\{x:\ \sup_n x_n\ge5\}$ --- борелевское.}
\end{task}
Осторожно: $\{\sup x_n\ge5\}\neq\bigcup_n\{x_n\ge5\}$~--- супремум может не
достигаться (например, $x_n=5-\frac1n$). Правильно: $\sup x_n\ge5$ $\iff$ для
любого $m$ найдётся член, больший $5-\frac1m$:
\[
\{\sup x_n\ge5\}=\bigcap_{m=1}^\infty\bigcup_{n=1}^\infty\{x_n>5-\tfrac1m\}.
\]
Счётные операции над цилиндрами~--- множество борелевское.

\begin{task}{Группа 2, б) момент первого попадания последовательности выше 2 --- случайная величина.}
\end{task}
Пусть $X_1,X_2,\dots$~--- случайные величины,
$\tau=\min\{n:\ X_n>2\}$ ($\tau=\infty$, если таких $n$ нет). Тогда
\[
\{\tau=n\}=\{X_1\le2\}\cap\dots\cap\{X_{n-1}\le2\}\cap\{X_n>2\}\in\F,
\qquad
\{\tau=\infty\}=\bigcap_{n}\{X_n\le2\}\in\F .
\]
Значит, $\tau$~--- случайная величина со значениями в $\N\cup\{\infty\}$.
Более того, $\{\tau=n\}\in\sigma(X_1,\dots,X_n)$, т.е. $\tau$~--- момент
остановки.

\begin{task}{Группа 3, а) множество последовательностей, сходящихся к 2, --- борелевское.}
\end{task}
Определение предела: $\forall\varepsilon>0\ \exists N\ \forall n\ge N:\
|x_n-2|<\varepsilon$. Заменяя $\varepsilon$ на $\frac1m$:
\[
\{x_n\to2\}=\bigcap_{m=1}^\infty\bigcup_{N=1}^\infty\bigcap_{n\ge N}\{|x_n-2|<\tfrac1m\},
\]
а $\{|x_n-2|<\frac1m\}=\{2-\frac1m<x_n<2+\frac1m\}$~--- цилиндр.

\begin{task}{Группа 3, б) момент первого троекратного повторения значения --- случайная величина.}
\end{task}
Понимаем так: $\tau$~--- первый момент $n$, когда значение $X_n$ встречается
в третий раз, т.е.
$\tau=\min\{n:\ \exists\,i<j<n:\ X_i=X_j=X_n\}$ ($\min\varnothing=\infty$).
Сначала: $\{X_i=X_j\}\in\F$, так как $X_i-X_j$~--- случайная величина и
$\{X_i=X_j\}=\{X_i-X_j=0\}$. Далее
\[
\{\tau\le n\}=\bigcup_{i<j<k\le n}\bigl(\{X_i=X_j\}\cap\{X_j=X_k\}\bigr)\in\F
\]
(конечное объединение событий), и
$\{\tau=n\}=\{\tau\le n\}\setminus\{\tau\le n-1\}\in\F$,
$\{\tau=\infty\}=\Omega\setminus\bigcup_n\{\tau\le n\}\in\F$.
Значит, $\tau$~--- случайная величина (и даже момент остановки).

\begin{zam}
Если понимать «троекратное повторение» как три \emph{подряд} одинаковых
значения, доказательство то же:
$\{\tau=n\}=\{X_{n-2}=X_{n-1}=X_n\}\setminus\bigcup_{k<n}\{X_{k-2}=X_{k-1}=X_k\}$.
\end{zam}

\section{Базовые задачи}\label{sec:basic}

\subsection{Задача: линейные комбинации определяют распределение}\label{z:cramer}

\begin{task}{Группа 1, №1. Показать, что зная распределения любой величины
$\lambda_1X_1+\dots+\lambda_kX_k$, мы однозначно можем восстановить распределение
последовательности $\{X_k\}$.}
\end{task}
\textbf{Шаг 1: восстанавливаем х.ф. каждого начального отрезка.} Зафиксируем
$k$ и $\lambda=(\lambda_1,\dots,\lambda_k)$. Величина
$Z_\lambda=\lambda_1X_1+\dots+\lambda_kX_k$ имеет известное распределение,
значит, известна её х.ф., в частности в точке $1$:
\[
\varphi_{Z_\lambda}(1)=\EE e^{iZ_\lambda}=\EE e^{i(\lambda_1X_1+\dots+\lambda_kX_k)}
=\varphi_{(X_1,\dots,X_k)}(\lambda_1,\dots,\lambda_k).
\]
Так как $\lambda$ произвольно, известна вся х.ф. вектора $(X_1,\dots,X_k)$.

\textbf{Шаг 2: к.м.р.} По теореме единственности х.ф. вектора определяет его
распределение. Значит, известны распределения всех векторов
$(X_1,\dots,X_k)$, $k\ge1$, а через них и все к.м.р. (распределение
$(X_{t_1},\dots,X_{t_k})$~--- проекция распределения
$(X_1,\dots,X_{\max t_j})$).

\textbf{Шаг 3: распределение последовательности.} К.м.р. задают значения
меры $\PP_X$ на всех цилиндрах; цилиндры~--- $\pi$-система, порождающая
$\B(\R^\infty)$. По теореме о $\pi$-системах $\PP_X$ определена однозначно.
\qed

\subsection{Задача о рабочем у конвейера}\label{z:worker}

\begin{task}{Группа 1, №2 (= Д2 группы 3). Каждую минуту поступает изделие,
с вероятностью $p$ оно дефектно. Проверка занимает полминуты, дефектное
изделие рабочий 5 минут исправляет (изделия копятся в очереди). $X_n$~--- число
изделий у рабочего в момент $n$, $Y_n$~--- время, которое он уже потратил на
текущее изделие. Является ли $X_n$ ЦМ? Пара $(X_n,Y_n)$?}
\end{task}

\paragraph{Уточнение модели.} Изделия приходят в моменты $0,1,2,\dots$ и
обрабатываются по очереди. Дефектность изделий~--- независимые индикаторы
$\varepsilon_1,\varepsilon_2,\dots$ с $\PP(\varepsilon=1)=p$ (нумеруем в порядке
обработки). Годное изделие занимает $\frac12$ минуты, дефектное~---
$\frac12+5=5\frac12$ минуты. $X_n$~--- число изделий у рабочего сразу после
прихода изделия в момент $n$ (включая то, что в работе); $Y_n$~--- сколько
времени он уже потратил на текущее изделие; закончив изделие, рабочий сразу
берёт следующее. Считаем $0<p<1$ (при $p=0$ или $p=1$ всё детерминировано, и
любая детерминированная последовательность тривиально марковская).

\paragraph{Ответ: $X_n$ не ЦМ, пара $(X_n,Y_n)$~--- ЦМ.}

\paragraph{Почему $X_n$ не ЦМ.} Посчитаем две траектории.

\emph{Траектория 1:} изделия $0,\dots,7$ годные, изделие $8$ дефектное. Каждое
годное изделие уходит за полминуты, поэтому $X_0=\dots=X_8=1$. Изделие $8$
проверяется на $[8;8{,}5]$ и ремонтируется на $[8{,}5;13{,}5]$, так что в моменты
$9$ и $10$ у рабочего $2$ и $3$ изделия: $X_9=2$, $X_{10}=3$ наверняка.

\emph{Траектория 2:} изделие $0$ дефектное, изделия $1,\dots,7$ годные.
Ремонт идёт на $[0{,}5;5{,}5]$, очередь растёт: $X_1=2,\dots,X_5=6$. Потом
годные изделия уходят по два за минуту, а приходит одно, и очередь убывает:
$X_6=5$, $X_7=4$, $X_8=3$, $X_9=2$. В момент $9$ рабочий берёт изделие $8$,
про которое ничего не известно: если оно годное, $X_{10}\le2$, если
дефектное~--- $X_{10}=3$.

Обе предыстории имеют положительную вероятность, и обе кончаются
в $X_9=2$, но
\begin{gather*}
\PP(X_{10}=3\mid X_0=\dots=X_8=1,\ X_9=2)=1,\\
\PP(X_{10}=3\mid X_0=1,X_1=2,\dots,X_5=6,X_6=5,\dots,X_8=3,\ X_9=2)=p<1 .
\end{gather*}
Условная вероятность будущего зависит не только от настоящего~--- $X_n$ не
ЦМ. Причина: по одному числу изделий не видно, идёт ли сейчас ремонт и сколько
ему осталось, а это прошлое «помнит».

\paragraph{Почему $(X_n,Y_n)$~--- ЦМ.} Заметим, что в целые моменты $Y_n$
принимает значения $0,\frac12,1,\dots,5$, причём $Y_n\ge\frac12$ бывает
только у изделия в ремонте (годное изделие к моменту, когда на него потрачено
$\frac12$, уже закончено). Поэтому пара $(X_n,Y_n)$ знает всё существенное о
текущем моменте: длину очереди, идёт ли ремонт и сколько ему осталось
($5\frac12-Y_n$). Дальнейшая эволюция на отрезке $[n,n+1]$~--- детерминированная
функция от $(X_n,Y_n)$ и индикаторов дефектности изделий, проверка которых
ещё не закончена к моменту $n$. Пусть $K_n$~--- число изделий, проверка
которых закончилась к моменту $n$; всё прошлое $(X_m,Y_m)_{m\le n}$~---
функция от $\varepsilon_1,\dots,\varepsilon_{K_n}$, причём $K_n$~--- момент
остановки для последовательности $\varepsilon$. Для последовательности
независимых одинаково распределённых величин «хвост» после момента остановки
$(\varepsilon_{K_n+1},\varepsilon_{K_n+2},\dots)$ не зависит от прошлого и
распределён так же, как вся последовательность. Итак,
$(X_{n+1},Y_{n+1})=F(X_n,Y_n,\text{свежий независимый шум})$, и по
лемме~\ref{lem:rec} пара~--- (однородная) ЦМ.

\subsection{Задача: к.м.р. неоднородной цепи и х.ф. суммы}\label{z:inhom}

\begin{task}{Группа 1, №3 (= №3 группы 3). а) Выразить
$\PP(X_{t_1}=i_1,\dots,X_{t_k}=i_k)$ через $p^0$ и м.в.п. $P_n$ неоднородной
ЦМ, $t_1<\dots<t_k$. б) Х.ф. $X_0+\dots+X_n$ через вектор $q^0$,
$q^0_j=p^0_je^{itj}$, и матрицу $Q$, $q_{lj}=p_{lj}e^{itj}$.}
\end{task}

Обозначения: $(P_m)_{ij}=\PP(X_{m+1}=j\mid X_m=i)$, и для $s\le t$
\[
P^{(s,t)}=P_sP_{s+1}\cdots P_{t-1}\qquad(P^{(s,s)}=E).
\]

\textbf{а)} Начнём с полной траектории до момента $T=t_k$:
\[
\PP(X_0=j_0,\dots,X_T=j_T)=p^0_{j_0}(P_0)_{j_0j_1}(P_1)_{j_1j_2}\cdots(P_{T-1})_{j_{T-1}j_T}.
\]
Искомая вероятность~--- сумма этих произведений по всем $j_m$ при
$m\notin\{t_1,\dots,t_k\}$, с фиксированными $j_{t_r}=i_r$. Сумма по
«свободным» индексам между $t_r$ и $t_{r+1}$~--- это ровно определение
матричного произведения:
$\sum_{j_{t_r+1},\dots,j_{t_{r+1}-1}}(P_{t_r})_{i_r\,j_{t_r+1}}\cdots(P_{t_{r+1}-1})_{j_{t_{r+1}-1}\,i_{r+1}}
=\bigl(P^{(t_r,t_{r+1})}\bigr)_{i_ri_{r+1}}$.
Аналогично сумма по $j_0,\dots,j_{t_1-1}$ даёт $(p^0P^{(0,t_1)})_{i_1}$. Итого
\[
\boxed{\PP(X_{t_1}=i_1,\dots,X_{t_k}=i_k)=
\bigl(p^0P_0\cdots P_{t_1-1}\bigr)_{i_1}\prod_{r=1}^{k-1}
\bigl(P_{t_r}P_{t_r+1}\cdots P_{t_{r+1}-1}\bigr)_{i_ri_{r+1}} .}
\]
Смысл: сначала надо дойти до $i_1$ за время $t_1$, потом из $i_1$ в $i_2$ за
время $t_2-t_1$ и т.д.; каждый этап~--- матрица переходов за несколько шагов.

\textbf{б)} Состояния~--- числа $j$ (значения цепи). Пусть $S=X_0+\dots+X_n$.
По формуле полной вероятности
\[
\varphi_S(t)=\EE e^{itS}=\sum_{j_0,\dots,j_n}\PP(X_0=j_0,\dots,X_n=j_n)\,e^{it(j_0+\dots+j_n)}
=\sum_{j_0,\dots,j_n}\bigl(p^0_{j_0}e^{itj_0}\bigr)\bigl(p_{j_0j_1}e^{itj_1}\bigr)\cdots\bigl(p_{j_{n-1}j_n}e^{itj_n}\bigr).
\]
Каждый множитель $e^{itj_m}$ «приклеили» к тому элементу, в котором $j_m$~---
номер \emph{столбца}. Это и есть элементы $q^0$ и $Q$:
\[
\varphi_S(t)=\sum_{j_0,\dots,j_n}q^0_{j_0}q_{j_0j_1}\cdots q_{j_{n-1}j_n}
=\boxed{\,q^0Q^n\mathbf 1\,},
\]
где $\mathbf 1=(1,\dots,1)^{\mathsf T}$ (сумма по последнему индексу $j_n$~---
это умножение на столбец из единиц). Для неоднородной цепи так же
$\varphi_S(t)=q^0Q_0Q_1\cdots Q_{n-1}\mathbf 1$, где
$(Q_m)_{lj}=(P_m)_{lj}e^{itj}$. (В условии опечатка: $q^0_j=p^0_je^{itj}$.)

\subsection{Задача: момент остановки}\label{z:stop}

\begin{task}{Группа 1, №4 (= Д1 группы 3). $X_n$~--- ЦМ, $N\ge1$ такая, что
$\{N=n\}\in\sigma(X_1,\dots,X_n)$. Показать, что
$\PP(X_{N+1}=j\mid X_N=i,\ X_l=x_l,\ l\le N)=p_{ij}$.}
\end{task}
Условие «$X_l=x_l,\ l\le N$» с данным набором $(x_1,\dots,x_n)$, $x_n=i$,~---
это событие
\[
B_n=\{N=n,\ X_1=x_1,\dots,X_n=x_n\}.
\]
\textbf{Ключевое наблюдение.} $\{N=n\}\in\sigma(X_1,\dots,X_n)$, а при дискретных
состояниях $\sigma(X_1,\dots,X_n)$ состоит из объединений «атомов»
$A=\{X_1=x_1,\dots,X_n=x_n\}$. Поэтому на атоме $A$ событие $\{N=n\}$ либо
выполнено целиком, либо не выполнено нигде. Если $\PP(B_n)>0$, то
$B_n=A$.

Тогда
\[
\PP(X_{N+1}=j\mid B_n)=\frac{\PP(X_{n+1}=j,\ X_1=x_1,\dots,X_n=x_n)}{\PP(X_1=x_1,\dots,X_n=x_n)}
=p_{x_nj}=p_{ij}
\]
по обычному марковскому свойству (на $B_n$ имеем $N=n$, так что
$X_{N+1}=X_{n+1}$). Если же условие~--- объединение нескольких таких $B_n$
(например, только $\{X_N=i\}$), то вероятность совместного события равна
$\sum_n p_{ij}\PP(B_n)=p_{ij}\PP(\bigsqcup B_n)$, и снова получаем $p_{ij}$.
\qed

Смысл: момент остановки не «подглядывает» в будущее, поэтому после него цепь
продолжает жить с той же м.в.п. (строго марковское свойство).

\subsection{Задача: функция от ЦМ, не являющаяся ЦМ}\label{z:func}

\begin{task}{Группа 1, Д1 (= №4 группы 3). Показать, что существуют ЦМ $X_n$ с
тремя состояниями и функция $h$, т.ч. $h(X_n)$ не ЦМ.}
\end{task}
Пусть $X_n$ детерминированно обходит цикл $1\to2\to3\to1$:
\[
P=\begin{pmatrix}0&1&0\\0&0&1\\1&0&0\end{pmatrix},\qquad p^0=\bigl(\tfrac13,\tfrac13,\tfrac13\bigr),
\]
и $h(1)=h(2)=a$, $h(3)=b$. Положим $Y_n=h(X_n)$. Тогда при $n\ge1$
\begin{itemize}[nosep]
\item $\{Y_{n-1}=a,\ Y_n=a\}=\{X_{n-1}=1,\ X_n=2\}$, следовательно $X_{n+1}=3$ и
$\PP(Y_{n+1}=a\mid Y_n=a,\ Y_{n-1}=a)=0$;
\item $\{Y_{n-1}=b,\ Y_n=a\}=\{X_{n-1}=3,\ X_n=1\}$, следовательно $X_{n+1}=2$ и
$\PP(Y_{n+1}=a\mid Y_n=a,\ Y_{n-1}=b)=1$.
\end{itemize}
Оба условия имеют вероятность $\frac13$. Настоящее одно ($Y_n=a$), а
прогноз разный~--- $Y_n$ не ЦМ. Склейка состояний $1$ и $2$ уничтожила
информацию о том, «в каком месте цикла» мы находимся.

\subsection{Задача о блуждании по многоугольнику}\label{z:polygon}

\begin{task}{Группа 1, Д2 (= Д1 группы 2). Является ли блуждание по вершинам
$A_1,\dots,A_n$ многоугольника ЦМ, если, стартуя из $A_1$, а) частица движется
по часовой; б) частица в начале случайно выбирает направление и движется в
нём; в) частица из $A_i$, $i\ne n$, переходит в соседние вершины с
вероятностями $p,1-p$, из $A_n$ возвращается туда, откуда пришла ($n$ чётно).}
\end{task}
Считаем $n\ge3$, по часовой~--- $A_i\to A_{i+1}$ (индексы по модулю $n$).

\textbf{а) Да.} $X_{k+1}=A_{i+1}$ при $X_k=A_i$~--- лемма~\ref{lem:rec} без шума.
М.в.п.: $p_{A_i,A_{i+1}}=1$, остальные нули.

\textbf{б) Нет.} Пусть направление выбирается с вероятностями
$\frac12,\frac12$. В момент $n$ частица в обоих случаях снова в $A_1$ (обошла
весь многоугольник), но предыстории разные:
\begin{itemize}[nosep]
\item по часовой: $X_{n-1}=A_n$, $X_n=A_1$, и далее $X_{n+1}=A_2$;
\item против часовой: $X_{n-1}=A_2$, $X_n=A_1$, и далее $X_{n+1}=A_n$.
\end{itemize}
Значит,
$\PP(X_{n+1}=A_2\mid X_n=A_1,\ X_{n-1}=A_n)=1$, а
$\PP(X_{n+1}=A_2\mid X_n=A_1)=\frac12$. Не ЦМ. (Пара «вершина, направление»
уже ЦМ.)

\textbf{в) Нет} (при $0<p<1$). Из $A_n$ частица идёт туда, откуда пришла, то
есть будущее зависит от предыдущей вершины. Нужно предъявить две предыстории
одной длины. Пусть частица из $A_i$ ($i\ne n$) идёт к $A_{i+1}$ с вероятностью
$p$ и к $A_{i-1}$ с вероятностью $1-p$.
\begin{itemize}[nosep]
\item Путь $\alpha$: $A_1\to A_2\to A_1\to A_2\to\dots\to A_1$ (всего $n-2$
шагов~--- чётное число, поэтому заканчиваем в $A_1$), затем $A_1\to A_n$. Длина
$n-1$, предпоследняя вершина $A_1$.
\item Путь $\beta$: $A_1\to A_2\to\dots\to A_{n-1}\to A_n$. Длина $n-1$,
предпоследняя вершина $A_{n-1}$.
\end{itemize}
Оба пути имеют положительную вероятность (здесь и нужна чётность $n$). Тогда
\[
\PP(X_n=A_1\mid X_{n-1}=A_n,\ X_{n-2}=A_1,\ \dots)=1,\qquad
\PP(X_n=A_1\mid X_{n-1}=A_n,\ X_{n-2}=A_{n-1},\ \dots)=0 .
\]
Не ЦМ. Зато пара $(X_{k-1},X_k)$~--- ЦМ (расширение состояния).

\subsection{Задача: условия согласованности}\label{z:consist}

\begin{task}{Группа 2, №1 и группа 3, №1. Сформулировать (обосновав
равносильность) условия согласованности в терминах а) характеристических
функций, б) функций распределения вектора $(X_1,\dots,X_k)$.}
\end{task}
Пусть $P_k$~--- меры на $\R^k$, $\varphi_k$~--- их х.ф., $F_k$~--- их ф.р.
Исходное условие согласованности:
\[
\text{(С)}\qquad P_{k+1}(B\times\R)=P_k(B),\qquad B\in\B(\R^k),\ k\ge1 .
\]
Левая часть~--- это мера $P_{k+1}\circ\pi^{-1}$, где $\pi(x_1,\dots,x_{k+1})=(x_1,\dots,x_k)$~---
проекция. То есть (С) означает: «проекция $P_{k+1}$ на первые $k$ координат
равна $P_k$».

\textbf{а) Через х.ф.:} $\varphi_{k+1}(\lambda_1,\dots,\lambda_k,0)=\varphi_k(\lambda_1,\dots,\lambda_k)$
для всех $\lambda\in\R^k$, $k\ge1$.

\emph{(С) $\Rightarrow$ (а).} По формуле замены переменной
\[
\varphi_{k+1}(\lambda,0)=\int_{\R^{k+1}}e^{i(\lambda_1x_1+\dots+\lambda_kx_k)}\,P_{k+1}(dx)
=\int_{\R^k}e^{i(\lambda,y)}\,(P_{k+1}\circ\pi^{-1})(dy)
=\int_{\R^k}e^{i(\lambda,y)}\,P_k(dy)=\varphi_k(\lambda).
\]
\emph{(а) $\Rightarrow$ (С).} Та же выкладка показывает, что
$\varphi_{k+1}(\lambda,0)$~--- х.ф. проекции $P_{k+1}\circ\pi^{-1}$. По условию
она совпадает с х.ф. $P_k$, а по теореме единственности совпадают и меры.

Если к.м.р. заданы для произвольных наборов индексов, добавляется условие
симметрии: $\varphi_{t_{\sigma(1)}\dots t_{\sigma(k)}}(\lambda_{\sigma(1)},\dots,\lambda_{\sigma(k)})
=\varphi_{t_1\dots t_k}(\lambda_1,\dots,\lambda_k)$ для любой перестановки
$\sigma$ (равносильность~--- так же, через единственность), а «зануление»
разрешается для любой координаты.

\textbf{б) Через ф.р.:}
$\lim\limits_{x_{k+1}\to+\infty}F_{k+1}(x_1,\dots,x_k,x_{k+1})=F_k(x_1,\dots,x_k)$.

\emph{(С) $\Rightarrow$ (б).} Множества
$(-\infty,x_1]\times\dots\times(-\infty,x_k]\times(-\infty,m]$ при
$m\to\infty$ возрастают к $(-\infty,x_1]\times\dots\times(-\infty,x_k]\times\R$.
По непрерывности меры
$F_{k+1}(x_1,\dots,x_k,m)\to P_{k+1}\bigl(\prod(-\infty,x_i]\times\R\bigr)=P_k\bigl(\prod(-\infty,x_i]\bigr)=F_k(x)$.
(Монотонность по $x_{k+1}$ позволяет перейти от $m\to\infty$ к
$x_{k+1}\to+\infty$.)

\emph{(б) $\Rightarrow$ (С).} Та же выкладка показывает, что предел слева~---
ф.р. проекции $P_{k+1}\circ\pi^{-1}$. Она совпадает с ф.р. $P_k$; ф.р.
определяет меру (множества $\prod(-\infty,x_i]$~--- $\pi$-система,
порождающая $\B(\R^k)$), значит, проекция равна $P_k$.

\subsection{Задача: подпоследовательности и функции от ЦМ}\label{z:subseq}

\begin{task}{Группа 2, №2. $X_n$~--- ЦМ. Будут ли $X_{2n}$, $X_n^3$,
$(X_n,X_{n+1})$ ЦМ?}
\end{task}
\textbf{$X_{2n}$~--- да.} Пусть $Z_n=X_{2n}$. По марковскому свойству
$X$ и уравнению Колмогорова -- Чепмена
\[
\PP(X_{2n+2}=j\mid X_{2n}=i,\ X_{2n-1}=\cdot,\dots,X_0=\cdot)=\bigl(P_{2n}P_{2n+1}\bigr)_{ij}
\]
--- при условии на \emph{всё} прошлое. Условие
$\{Z_n=i,Z_{n-1}=i_{n-1},\dots,Z_0=i_0\}$~--- объединение таких событий по
значениям $X$ в нечётные моменты; для каждого слагаемого условная вероятность
одна и та же, значит, и для объединения тоже $(P_{2n}P_{2n+1})_{ij}$. ЦМ с
м.в.п. $P_{2n}P_{2n+1}$ (в однородном случае $P^2$).

\textbf{$X_n^3$~--- да.} Функция $x\mapsto x^3$ инъективна на $\R$, поэтому
$\{X_n^3=i^3\}=\{X_n=i\}$ и все условные вероятности те же:
$\PP(X_{n+1}^3=j^3\mid X_n^3=i^3,\dots)=p_{ij}$.

\textbf{$(X_n,X_{n+1})$~--- да.} Обозначим $W_n=(X_n,X_{n+1})$. Тогда
\[
\PP\bigl(W_{n+1}=(j',k)\mid W_n=(i,j),\ W_{n-1},\dots\bigr)
=\PP(X_{n+2}=k,\ X_{n+1}=j'\mid X_{n+1}=j,\ X_n=i,\dots)=\1\{j'=j\}\,p_{jk}(n+1),
\]
что зависит только от текущего $W_n$. ЦМ со «ступенчатой» м.в.п.

\subsection{Задача: к.м.р. однородной ЦМ}\label{z:hom}

\begin{task}{Группа 2, №3. Выразить $\PP(X_{t_1}=i_1,\dots,X_{t_k}=i_k)$ через
$p^0$ и м.в.п. $P$ однородной ЦМ.}
\end{task}
Частный случай задачи~\ref{z:inhom} с $P_m\equiv P$, когда
$P_{t_r}\cdots P_{t_{r+1}-1}=P^{t_{r+1}-t_r}$:
\[
\boxed{\PP(X_{t_1}=i_1,\dots,X_{t_k}=i_k)=(p^0P^{t_1})_{i_1}\,(P^{t_2-t_1})_{i_1i_2}\cdots(P^{t_k-t_{k-1}})_{i_{k-1}i_k}.}
\]
Вывод тот же: записать вероятность полной траектории до $t_k$ и просуммировать
по состояниям в промежуточные моменты; суммы превращаются в степени матрицы
$P$ (уравнение Колмогорова -- Чепмена). Другая запись того же вывода~--- по
формуле умножения и марковскому свойству:
$\PP(X_{t_{r+1}}=i_{r+1}\mid X_{t_r}=i_r,\dots,X_{t_1}=i_1)=(P^{t_{r+1}-t_r})_{i_ri_{r+1}}$.

\subsection{Задача: стационарность и обращение времени}\label{z:station}

\begin{task}{Группа 2, №4. $X_n$~--- однородная ЦМ. а) Показать, что $X_n\eqd X_m$
тогда и только тогда, когда начальное распределение~--- левый собственный
вектор м.в.п. с собственным значением 1. б) При этом условии показать, что
$Y_n=X_{N-n}$, $n\le N$, тоже ЦМ. в) Когда $Y_n$ имеет ту же м.в.п., что и $X_n$?
Ответ в терминах собственного вектора $P^{\mathsf T}$.}
\end{task}

\textbf{а)} Понимаем «$X_n\eqd X_m$ для всех $n,m$». Распределение $X_n$~---
вектор $p^n=p^0P^n$.

($\Leftarrow$) Если $p^0P=p^0$, то по индукции
$p^0P^n=(p^0P^{n-1})P=p^0P=p^0$ для всех $n$: все $X_n$ распределены как $X_0$.

($\Rightarrow$) Если все $X_n$ одинаково распределены, в частности
$X_1\eqd X_0$, то $p^0P=p^1=p^0$.

Более того, тогда цепь стационарна в узком смысле:
$\PP(X_n=i_0,\dots,X_{n+k}=i_k)=p^n_{i_0}p_{i_0i_1}\cdots p_{i_{k-1}i_k}$
не зависит от $n$. (Если $X_n\eqd X_m$ лишь для одной пары $n\ne m$, это
слабее: например, у периодической цепи с $P=\begin{pmatrix}0&1\\1&0\end{pmatrix}$
$X_0\eqd X_2$ при любом $p^0$.)

\textbf{б)} Обозначим $\pi=p^0$ (инвариантное распределение), считаем $\pi_i>0$.
Для $n<N$
\[
\PP(Y_{n+1}=j\mid Y_n=i,Y_{n-1}=i_{n-1},\dots,Y_0=i_0)
=\frac{\PP(X_{N-n-1}=j,\ X_{N-n}=i,\ X_{N-n+1}=i_{n-1},\dots,X_N=i_0)}
{\PP(X_{N-n}=i,\ X_{N-n+1}=i_{n-1},\dots,X_N=i_0)} .
\]
По формуле совместного распределения и стационарности
($\PP(X_m=\cdot)=\pi$):
\[
=\frac{\pi_j\,p_{ji}\,p_{i\,i_{n-1}}\cdots p_{i_1i_0}}{\pi_i\,p_{i\,i_{n-1}}\cdots p_{i_1i_0}}
=\frac{\pi_jp_{ji}}{\pi_i}.
\]
Ответ зависит только от $i$ и $j$, так что $Y_n$~--- (однородная) ЦМ с м.в.п.
\[
q_{ij}=\frac{\pi_jp_{ji}}{\pi_i}.
\]
(Проверка стохастичности:
$\sum_jq_{ij}=\frac{1}{\pi_i}\sum_j\pi_jp_{ji}=\frac{(\pi P)_i}{\pi_i}=1$~---
здесь как раз используется $\pi P=\pi$. Без стационарности обращённая цепь
тоже марковская, но неоднородная.)

\textbf{в)} $q_{ij}=p_{ij}$ $\iff$
\[
\boxed{\pi_ip_{ij}=\pi_jp_{ji}\quad\text{для всех }i,j,}
\]
где $\pi^{\mathsf T}$~--- собственный вектор $P^{\mathsf T}$ с собственным
значением $1$ ($P^{\mathsf T}\pi^{\mathsf T}=\pi^{\mathsf T}$~--- то же, что $\pi P=\pi$),
нормированный на сумму $1$. В матричной форме: матрица $D_\pi P$ симметрична,
где $D_\pi=\mathrm{diag}(\pi_1,\pi_2,\dots)$. Это условие детального баланса
(обратимость цепи): поток вероятности из $i$ в $j$ равен потоку из $j$ в $i$.

\subsection{Задача о бросаниях кубика}\label{z:dice}

\begin{task}{Группа 2, Д2. Кто из них является ЦМ? а) $X_n$~--- максимум чисел на
кубике при первых $n$ бросках; б) число шестёрок при первых $n$ бросках;
в) $B_n$~--- время с момента $n$ до следующей шестёрки. Для всех ЦМ найти м.в.п.}
\end{task}
Пусть $\xi_1,\xi_2,\dots$~--- независимые результаты бросков, равномерные на
$\{1,\dots,6\}$.

\textbf{а) Да.} $X_{n+1}=\max(X_n,\xi_{n+1})$~--- лемма~\ref{lem:rec}.
Состояния $1,\dots,6$:
\[
p_{ij}=\begin{cases}0,&j<i,\\ i/6,&j=i\ (\text{выпало не больше }i),\\ 1/6,&j>i.\end{cases}
\qquad
P=\frac16\begin{pmatrix}
1&1&1&1&1&1\\0&2&1&1&1&1\\0&0&3&1&1&1\\0&0&0&4&1&1\\0&0&0&0&5&1\\0&0&0&0&0&6
\end{pmatrix}.
\]
Начальное распределение $X_1$~--- равномерное (или $X_0=0$ с $p_{0j}=\frac16$).

\textbf{б) Да.} $S_{n+1}=S_n+\1\{\xi_{n+1}=6\}$. Состояния $0,1,2,\dots$:
$p_{k,k+1}=\frac16$, $p_{k,k}=\frac56$, остальные нули; $S_0=0$.

\textbf{в) Да.} $B_n=\min\{k\ge1:\ \xi_{n+k}=6\}$. Если $B_n\ge2$, то
следующая шестёрка та же, и $B_{n+1}=B_n-1$. Если $B_n=1$, то $\xi_{n+1}=6$, и
$B_{n+1}$~--- время ожидания новой шестёрки по броскам $\xi_{n+2},\xi_{n+3},\dots$,
геометрическое: $\PP(B_{n+1}=m)=\left(\frac56\right)^{m-1}\frac16$.
При этом на событии $\{B_n=1\}$ вся предыстория $B_0,\dots,B_n$~--- функция от
$\xi_1,\dots,\xi_{n+1}$ (все «следующие шестёрки» для моментов $\le n$
наступают не позже $n+1$), а $B_{n+1}$ зависит только от $\xi_{n+2},\dots$ ---
независимость есть. Итого однородная ЦМ на $\{1,2,\dots\}$:
\[
p_{k,k-1}=1\ (k\ge2),\qquad p_{1,m}=\left(\tfrac56\right)^{m-1}\tfrac16\ (m\ge1),
\]
остальные нули. Начальное распределение: $B_0$ геометрическое с параметром
$\frac16$. (Это цепь «остаточного времени ожидания».)

\subsection{Задача о максимуме случайного блуждания}\label{z:walk}

\begin{task}{Группа 3, №2. $S_n$~--- простое симметричное случайное блуждание,
$M_n=\max(S_0,\dots,S_n)$. Является ли ЦМ а) пара $(M_n,S_n)$? б) $M_n-S_n$?
в) само $M_n$?}
\end{task}
$S_0=0$, $S_{n+1}=S_n+\xi_{n+1}$, $\xi_k$ независимы,
$\PP(\xi=\pm1)=\frac12$. Ключевое соотношение:
$M_{n+1}=\max(M_n,S_{n+1})=\max(M_n,S_n+\xi_{n+1})$.

\textbf{а) Да.} $(M_{n+1},S_{n+1})=\bigl(\max(M_n,S_n+\xi_{n+1}),\ S_n+\xi_{n+1}\bigr)$
--- функция от $(M_n,S_n)$ и независимого $\xi_{n+1}$ (лемма~\ref{lem:rec}).
Переходы из $(m,s)$, $s\le m$:
\[
(m,s)\to(m,s-1)\ \text{с вер. }\tfrac12,\qquad
(m,s)\to(\max(m,s+1),\ s+1)\ \text{с вер. }\tfrac12 .
\]

\textbf{б) Да.} Пусть $R_n=M_n-S_n\ge0$. Тогда
\[
R_{n+1}=\max(M_n,S_n+\xi_{n+1})-S_n-\xi_{n+1}=\max(R_n-\xi_{n+1},\,0),
\]
снова лемма~\ref{lem:rec}. М.в.п. (блуждание с отражением в нуле):
$p_{r,r+1}=p_{r,r-1}=\frac12$ при $r\ge1$; $p_{0,1}=\frac12$ (шаг вниз
$\xi=-1$), $p_{0,0}=\frac12$ (шаг вверх: новый максимум).

\textbf{в) Нет.} Вероятность того, что максимум вырастет, зависит от того,
находится ли $S_n$ на максимуме, а по одному $M_n$ этого не видно. Явно, при
$n=3$ с $M_3=1$ (каждый путь из трёх шагов имеет вероятность $\frac18$):
\begin{itemize}[nosep]
\item история $(M_0,\dots,M_3)=(0,0,0,1)$ даётся только путём $(-,+,+)$,
$S_3=1=M_3$, поэтому $\PP(M_4=2\mid\text{эта история})=\frac12$;
\item история $(0,1,1,1)$ даётся путями $(+,-,+)$ ($S_3=1$) и
$(+,-,-)$ ($S_3=-1$), поэтому
$\PP(M_4=2\mid\text{эта история})=\frac{\frac18\cdot\frac12+\frac18\cdot0}{\frac28}=\frac14$.
\end{itemize}
Настоящее одно ($M_3=1$), а вероятности разные: $M_n$ не ЦМ. Пара
$(M_n,S_n)$ из пункта а)~--- как раз то расширение состояния, которое
восстанавливает марковость.

\section{Факультативные задачи}\label{sec:fac}

\subsection*{Вспомогательный факт: где взять независимые величины на $[0,1]$}

На $([0,1],\B,\lambda)$ двоичные цифры $\omega=0.\omega_1\omega_2\dots$ числа
$\omega$~--- независимые величины с $\PP(\omega_k=0)=\PP(\omega_k=1)=\frac12$
(событие «первые $k$ цифр заданы»~--- отрезок длины $2^{-k}$). Разобьём $\N$ на
счётное число бесконечных непересекающихся множеств, например
$I_r=\{2^{r-1}(2m-1):\ m\ge1\}$, $r=1,2,\dots$, и положим
\[
U_r=\sum_{m\ge1}2^{-m}\omega_{2^{r-1}(2m-1)} .
\]
$U_r$ строятся по непересекающимся группам независимых цифр, поэтому
независимы, и каждая равномерна на $[0,1]$ (её двоичные цифры~--- снова
независимые честные монетки). Итак, на $([0,1],\B,\lambda)$ существует
последовательность независимых $U[0,1]$ величин $U_1,U_2,\dots$, а из них~---
величины с любыми распределениями: $F^{-1}(U_r)$ имеет ф.р. $F$.

\subsection{Ф-1.1. Цепь с двумя состояниями на $[0,1]$}

\begin{task}{Построить на $([0,1],\B([0,1]),\lambda)$ однородную ЦМ с двумя
состояниями, $p_{11}=p_{22}=p$ и начальным распределением $(0,1)$.}
\end{task}
Состояния $\{1,2\}$, $X_0\equiv2$. С каждым шагом цепь остаётся на месте с
вероятностью $p$ и переходит в другое состояние с вероятностью $1-p$.

\textbf{Способ 1 (через $U_r$).} $X_0=2$,
$X_{n+1}=X_n$ при $U_{n+1}\le p$ и $X_{n+1}=3-X_n$ при $U_{n+1}>p$. Это
$X_{n+1}=f(X_n,U_{n+1})$ с независимыми $U$, так что по лемме~\ref{lem:rec}
получаем однородную ЦМ с $p_{ii}=\PP(U\le p)=p$.

\textbf{Способ 2 (явное деление отрезка).} Положим $\omega^{(0)}=\omega$ и
\[
\omega^{(k+1)}=\begin{cases}\omega^{(k)}/p,&\omega^{(k)}<p\ \ (\text{«остаёмся»}),\\
(\omega^{(k)}-p)/(1-p),&\omega^{(k)}\ge p\ \ (\text{«переходим»}),\end{cases}
\]
$X_0=2$, $X_{k+1}=X_k$, если $\omega^{(k)}<p$, иначе $X_{k+1}=3-X_k$. По
индукции: множество $\omega$ с заданными первыми $k$ решениями
«остаться/перейти» есть отрезок длины $p^a(1-p)^{k-a}$ ($a$~--- число
«остаться»), и на нём $\omega^{(k)}$ равномерно пробегает $[0,1)$. Отсюда
$\PP(X_0=2,X_1=i_1,\dots,X_k=i_k)=\prod p_{i_{m}i_{m+1}}$ с нужной м.в.п. Геометрически:
$[0,1]$ делится на части длин $p$ и $1-p$, каждая часть~--- снова в той же
пропорции, и т.д.

\subsection{Ф-1.2. Гауссовская последовательность с $\mathrm{cov}(X_n,X_{n+k})=e^{-k}$}

\begin{task}{Построить на том же пространстве гауссовскую последовательность с
$\EE X_n=0$, $\mathrm{cov}(X_n,X_{n+k})=e^{-k}$.}
\end{task}
По подсказке ищем $X_{n+1}=aX_n+\xi_n$, $\xi_n\sim\mathcal N(0,\sigma^2)$
независимы и не зависят от $X_0,\dots,X_n$.

\emph{Подбор параметров.} $\mathrm{Var}\,X_n=\mathrm{cov}(X_n,X_n)=e^0=1$.
Для $k\ge1$: $X_{n+k}=a^kX_n+(\text{комбинация }\xi_n,\dots,\xi_{n+k-1})$, и
шумы не коррелируют с $X_n$, поэтому
$\mathrm{cov}(X_n,X_{n+k})=a^k\,\mathrm{Var}\,X_n=a^k$. Нужно $a^k=e^{-k}$,
т.е. $a=e^{-1}$. Дисперсия сохраняется: $1=a^2\cdot1+\sigma^2$, т.е.
$\sigma^2=1-e^{-2}$.

\emph{Построение.} Возьмём независимые $U_0,U_1,\dots\sim U[0,1]$ (см.
вспомогательный факт), $Z_r=\Phi^{-1}(U_r)\sim\mathcal N(0,1)$, где
$\Phi$~--- ф.р. $\mathcal N(0,1)$ (на множестве меры нуль $\{U_r\in\{0,1\}\}$
положим $Z_r=0$). Определим
\[
X_0=Z_0,\qquad X_{n+1}=e^{-1}X_n+\sqrt{1-e^{-2}}\,Z_{n+1}.
\]
\emph{Проверка.} Каждый вектор $(X_0,\dots,X_n)$~--- линейное преобразование
гауссовского вектора $(Z_0,\dots,Z_n)$ с независимыми координатами, значит,
гауссовский. $\EE X_n=0$. По индукции $\mathrm{Var}\,X_n=1$, а
$\mathrm{cov}(X_n,X_{n+k})=e^{-k}$, как посчитано выше. Гауссовская
последовательность определяется средними и ковариациями, так что построено ровно
то, что нужно.

\subsection{Ф-1.3. Процесс авторегрессии}

\begin{task}{$X_n=\sum_{k=0}^\infty a^k\xi_{n-k}$, $0<a<1$, $\xi_k$~--- н.о.р.
$\mathcal N(0,1)$. Показать, что это ЦМ, и найти переходный оператор.}
\end{task}
\emph{Корректность.} $\sum_ka^{2k}<\infty$, поэтому ряд сходится в среднем
квадратичном (и п.н.: $\EE\sum_ka^k|\xi_{n-k}|=\EE|\xi|\sum_ka^k<\infty$).
$X_n\sim\mathcal N\bigl(0,\sum a^{2k}\bigr)=\mathcal N\bigl(0,\frac1{1-a^2}\bigr)$.

\emph{Рекуррентное соотношение.} Выделим в $X_{n+1}$ первое слагаемое:
\[
X_{n+1}=\xi_{n+1}+\sum_{k\ge1}a^k\xi_{n+1-k}=\xi_{n+1}+a\sum_{k\ge0}a^k\xi_{n-k}=aX_n+\xi_{n+1}.
\]
Все $X_m$ при $m\le n$~--- функции от $\{\xi_l:\ l\le n\}$, а $\xi_{n+1}$ от них не
зависит. По лемме~\ref{lem:rec} $X_n$~--- однородная ЦМ.

\emph{Переходные характеристики.} При $X_n=x$ имеем
$X_{n+1}=ax+\xi_{n+1}\sim\mathcal N(ax,1)$:
\[
P(x,B)=\int_B\frac{1}{\sqrt{2\pi}}e^{-(y-ax)^2/2}\,dy,\qquad
f(y\mid x)=\frac{1}{\sqrt{2\pi}}e^{-(y-ax)^2/2},
\]
\[
\boxed{(Ph)(x)=\EE\,h(ax+\xi)=\int_\R h(ax+z)\,\frac{e^{-z^2/2}}{\sqrt{2\pi}}\,dz .}
\]
Инвариантное распределение~--- $\mathcal N\bigl(0,\frac1{1-a^2}\bigr)$, и
процесс стационарен (при $a=e^{-1}$ и нормировке шума это ровно процесс из Ф-1.2).

\subsection{Ф-1.4. Стационарность цепи с переходной плотностью}

\begin{task}{Когда цепь с переходной плотностью $f(y\mid x)$ стационарна в узком
смысле? Привести пример начального распределения для
$f(y\mid x)=2(y-x)$ при $0<x<y<1$, $f(y\mid x)=2(1-x+y)$ при $0<y<x<1$.}
\end{task}
\emph{Общее условие.} Совместная плотность
$(X_n,\dots,X_{n+k})$ равна $p_n(x_0)f(x_1\mid x_0)\cdots f(x_k\mid x_{k-1})$,
где $p_n$~--- плотность $X_n$. Она не зависит от $n$ $\iff$ $p_n=p_0$ для всех
$n$ $\iff$ $p_1=p_0$, т.е.
\[
\boxed{p_0(y)=\int f(y\mid x)\,p_0(x)\,dx\quad\text{для п.в. }y}
\]
(начальная плотность инвариантна; дальше по индукции).

\emph{Проверка, что $f$~--- переходная плотность.} При фиксированном $x$
\[
\int_x^12(y-x)\,dy+\int_0^x2(1-x+y)\,dy=(1-x)^2+\bigl(2x(1-x)+x^2\bigr)=1 .
\]
\emph{Пример: $p_0\equiv1$ на $(0,1)$.} Проверяем инвариантность:
\[
\int_0^1f(y\mid x)\,dx=\int_0^y2(y-x)\,dx+\int_y^12(1-x+y)\,dx
=y^2+\bigl[2(1+y)(1-y)-(1-y^2)\bigr]=y^2+(1-y^2)=1 .
\]
Значит, равномерное распределение на $(0,1)$ инвариантно, и цепь с $X_0\sim U(0,1)$
стационарна в узком смысле.

\emph{Почему так получилось.} $f(y\mid x)=g\bigl((y-x)\bmod1\bigr)$ с
$g(u)=2u$ на $(0,1)$: $X_{n+1}=(X_n+\eta_{n+1})\bmod1$, где $\eta$ независимы
с плотностью $2u$. Это случайный поворот окружности, а равномерная мера на
окружности не меняется при поворотах.

\subsection{Ф-1.5. Порядковые статистики}

\begin{task}{$\xi_i$, $i\le m$~--- н.о.р. $U[0,1]$, $X_n=\xi_{(n)}$ ($n$-я
порядковая статистика). Показать, что $X_1,\dots,X_m$~--- неоднородная ЦМ, и
найти её переходную плотность.}
\end{task}
\emph{Совместная плотность.} Вектор всех порядковых статистик
$(X_1,\dots,X_m)$ имеет плотность $m!$ на множестве
$\{0<x_1<\dots<x_m<1\}$ (каждый из $m!$ порядков равновероятен, а плотность
$(\xi_1,\dots,\xi_m)$ равна $1$ на кубе).

\emph{Плотность начального отрезка.} Интегрируем по $x_{n+1},\dots,x_m$:
\[
\int_{x_n<x_{n+1}<\dots<x_m<1}dx_{n+1}\cdots dx_m=\frac{(1-x_n)^{m-n}}{(m-n)!}
\]
(объём симплекса в кубе $(x_n,1)^{m-n}$). Поэтому
\[
g_n(x_1,\dots,x_n)=\frac{m!}{(m-n)!}(1-x_n)^{m-n},\qquad0<x_1<\dots<x_n<1 .
\]
\emph{Условная плотность.}
\[
f(x_{n+1}\mid x_1,\dots,x_n)=\frac{g_{n+1}}{g_n}
=\frac{(m-n)(1-x_{n+1})^{m-n-1}}{(1-x_n)^{m-n}},\qquad x_n<x_{n+1}<1 .
\]
Она зависит только от $x_n$ (и от $n$), значит, $X_n$~--- ЦМ; зависимость от
$n$ означает, что цепь неоднородная. Переходная плотность
\[
\boxed{f_n(y\mid x)=\frac{(m-n)(1-y)^{m-n-1}}{(1-x)^{m-n}},\qquad0<x<y<1,\ \ n=1,\dots,m-1,}
\]
а начальная плотность $X_1=\min\xi_i$ равна $m(1-x)^{m-1}$.

\emph{Смысл.} При $X_n=x$ оставшиеся $m-n$ точек~--- независимые равномерные
на $(x,1)$, а $X_{n+1}$~--- их минимум:
$\PP(X_{n+1}>y\mid X_n=x)=\bigl(\frac{1-y}{1-x}\bigr)^{m-n}$; производная по $y$
со знаком минус даёт ту же формулу.

\end{document}
